\origpage{55}{GRAPH MINORS. III}
\origfig{figures/fig3}{135bp}{\textsc{Figure} 3}

\stmt
(2.3) \textit{Let $G$ be a simple planar graph, with a drawing $M$ such that $\rho(M) \le d$ and $M$ has a $d$-shell $C_1, \dots, C_d$. Then $G$ is isomorphic to a minor of $N(d, \lambda)$ for some sufficiently large $\lambda$.}

\textit{Proof.} We observe that the graph obtained from $G$ by deleting $V(C_1), \dots, V(C_d)$ is a forest, since if the corresponding drawing had a finite region $R$, we would have $d(R) > d$ in $M$ which is impossible. But any forest is isomorphic to a minor of one of $Y_1$, $Y_2$, $Y_3, \dots$. The rest of the proof is obvious, although tedious to write out, and we omit it.

\fix{We give the omitted part. If $d = 0$, then $G$ itself is a forest, and so it is isomorphic to a minor of some $Y_\lambda = N(0, \lambda)$. So we assume that $d \ge 1$. Let $F$ be the forest above. By (S1) its vertices are inside $C_d$, and so all of $F$ is inside $C_d$, since its edges do not meet $C_d$.

\textit{Step 1: we may assume that $F$ is a tree or is null.} Suppose that $F$ has at least two components, let $F_1$ be one of them, and let $F_2$ be the union of the others. We claim that the closure of some region of $M$ inside $C_d$ contains a vertex of $F_1$ and a vertex of $F_2$. Since $F$ has no circuit and is disjoint from $C_d$, the only circuit of the drawing formed by $F$ and $C_d$ is $C_d$. Two points not in a drawing in the plane that are not separated by any circuit of the drawing lie in the same region of it. So the set $\Omega$ of points inside $C_d$ that are not in $U(F)$ is a region of this drawing, and in particular it is connected. Every vertex inside $C_d$ is in $F$, and by (S2) every edge with its interior points inside $C_d$ is in $F$ or joins $C_d$ and $F$. So $\Omega$ is the union of the regions of $M$ inside $C_d$ and the edges of $M$ joining $C_d$ and $F$. Each such edge is in the closure of the two regions beside it. Let $v_1 \in V(F_1)$ and $v_2 \in V(F_2)$. Since $U(F)$ has no interior point, $v_1$ is in the closure of $\Omega$, and so in the closure of some region of $M$ inside $C_d$, and the same holds for $v_2$. Let $Z$ be the set of regions of $M$ inside $C_d$ whose closure contains a vertex of $F_1$. If a region in $Z$ has $v_2$ in its closure, the claim holds. Otherwise $Z$ is non-empty, and so is the set of the other regions of $M$ inside $C_d$. Suppose that no edge of $M$ in $\Omega$ has a region in $Z$ on one side and a region not in $Z$ on the other. Then the regions in $Z$, with the edges in $\Omega$ beside them, form an open subset of $\Omega$, the other regions with the edges beside them form another, and these two sets are disjoint, non-empty and cover $\Omega$, which is impossible since $\Omega$ is connected. So some edge $g$ in $\Omega$ has a region in $Z$ on one side and a region not in $Z$ on the other. It joins a vertex of $C_d$ and a vertex $w$ of $F$, and $w$ is in the closure of both regions. Since one of these regions is not in $Z$, $w$ is not in $F_1$, and so $w \in V(F_2)$. Thus the region beside $g$ in $Z$ satisfies the claim. Now let $q$ be a region as in the claim, with $u_1 \in V(F_1)$ and $u_2 \in V(F_2)$ in its closure, and add a new edge joining $u_1$ and $u_2$, drawn in $q$. The new graph is simple, since an edge of $G$ joining $u_1$ and $u_2$ would be an edge of $F$ joining two of its components. The circuits $C_1, \dots, C_d$ are still nested, and (S1) and (S2) still hold, since the new edge joins two vertices inside $C_d$. The new forest $F$ has one component fewer, and $G$ is isomorphic to a minor of the new graph. From now on we use only that $G$ is simple, that $C_1, \dots, C_d$ is a nested sequence of circuits satisfying (S1) and (S2), and that $F$ is a forest inside $C_d$, and these properties hold for the new graph. Repeating this step, we may assume that $F$ is a tree or is null.

\textit{Step 2: the edges.} Every edge of $G$ is in some $C_j$, or joins $C_j$ and $C_{j+1}$ for some $j < d$, or joins $C_d$ and $F$, or is in $F$. For an edge joining two vertices of $C_j$ is in $C_j$ by (S2), and an edge joining $C_j$ and $C_{j'}$ with $j' > j + 1$, or $C_j$ and $F$ with $j < d$, would have one end inside $C_{j+1}$ and the other outside, and so it would cross $C_{j+1}$. For $1 \le j \le d$, let $G_j$ be the subgraph of $G$ formed by $C_j, \dots, C_d$, $F$, and the edges between them. By (S1), $G_1 = G$.

\textit{Step 3: the cyclic orders.} We use the following consequence of the Jordan curve theorem.

\hspace*{12pt}(A) \textit{Let $D'$ be a closed disc in the interior of a closed disc $D$ in the plane, and let $A_1, \dots, A_q$ be arcs, each with one end in $bd(D')$, the other end in $bd(D)$, and its other points in the interior of $D$ and not in $D'$, such that no point other than an end is in two of them. Then, going once clockwise around $bd(D')$ from any point that is not an end of one of them, we meet their ends in some order $A_1, \dots, A_q$, and going once clockwise around $bd(D)$ from a suitable point, we meet their ends in the same order, where arcs with a common end are met together.}

For two of the arcs, together with a portion of $bd(D')$ and a portion of $bd(D)$, bound a closed disc that lies between $bd(D')$ and $bd(D)$, such that both portions run clockwise from one of the two arcs to the other. Every other arc without its ends is connected and misses these two arcs, and so it lies in this disc or outside it. So any three of the arcs meet both circles in the same cyclic order, and (A) follows.

Let $j < d$. Each edge between $C_j$ and $C_{j+1}$ has its interior points between $C_j$ and $C_{j+1}$, since its interior is connected, misses both circuits, and has an end on each. By (A), applied to the discs bounded by $C_{j+1}$ and $C_j$, these edges can be listed so that their ends in $C_{j+1}$, and also their ends in $C_j$, are in the clockwise cyclic order of these circuits.

Now suppose that $F$ is a tree. Choose a closed disc $K$ inside $C_d$ that contains $U(F)$ in its interior and lies so close to $U(F)$ that it meets no edge of $M$ other than those of $F$ and those joining $F$ and $C_d$, and meets each edge $g$ joining $F$ and $C_d$ only in points near the end of $g$ in $F$. (Such discs exist since $U(F)$ is a tree. By a homeomorphism of the plane we may assume that the edges of $F$ are polygonal, and then the points within a small enough distance of $U(F)$ form one.) Go once clockwise around $bd(K)$, starting from a point near a vertex $t$ of $F$ that is on no edge of $M$. On the way we pass along the edges of $F$, each twice, and this gives a closed walk $W = w_0, w_1, \dots, w_m$ of $F$ with $w_0 = w_m = t$ which passes each edge of $F$ twice. Each $k$ with $0 \le k \le m$ is a \textit{visit} of $W$ to $w_k$, and there $bd(K)$ passes near $w_k$, through a sector at $w_k$ between two consecutive edges of $F$ (for $k = 0$ and $k = m$, through the parts of such a sector after and before the starting point). An edge $g$ joining $w_k$ and $C_d$ lies near $w_k$ in one of these sectors, and we say that $g$ \textit{leaves $F$ at} the corresponding visit. The part of $g$ after its last point in $K$ is an arc from $bd(K)$ to $C_d$, and its end in $bd(K)$ is where $bd(K)$ passes this visit. By (A), applied to $K$ and the disc bounded by $C_d$, the edges joining $F$ and $C_d$ can be listed so that the visits at which they leave $F$ are in the order of $W$, starting from visit $0$, and their ends in $C_d$ are in the clockwise cyclic order of $C_d$.

\textit{Step 4: two facts about $N(d', \lambda)$.} First, for $d' \ge 1$, $N(d', \lambda)$ is isomorphic to a minor of $N(d', \lambda + 1)$, in such a way that each vertex of a circle is represented by a set containing two consecutive vertices of the corresponding circle. To see this, number the radial lines of $N(d', \lambda + 1)$ as $1, \dots, 3 \cdot 2^\lambda$ in order, so that for each $i$ the end-vertices of $Y_{\lambda + 1}$ on the lines $2i - 1$ and $2i$ have a common neighbour $u_i$. On each circle, contract the edges between the lines $2i - 1$ and $2i$, and delete the radial edges on the even lines, which are now parallel to those on the odd lines. Then contract the edge joining each $u_i$ to the end-vertex on the line $2i - 1$, and delete the edge joining it to the end-vertex on the line $2i$. Deleting the end-vertices of $Y_{\lambda + 1}$ leaves $Y_\lambda$, with the vertices $u_i$ as its end-vertices, in the same order. So the result is $N(d', \lambda)$, as required. We call a model of $G_j$ as a minor of $N(d - j + 1, \lambda)$ \textit{cyclic} if each vertex of $C_j$ is represented by a set that contains an interval of the outer circle, and these intervals cover that circle and are in the clockwise cyclic order of $C_j$. Composing a cyclic model with the minor above doubles the length of each of its intervals, and so the intervals of a cyclic model may be made as long as we wish.

Second, let $B_h$ $(h \ge 1)$ be the rooted binary tree in which each vertex at distance less than $h - 1$ from the root has two children, a left one and a right one, and all leaves are at distance $h - 1$ from the root. Its leaves are ordered from left to right. We show by induction on $|E(T)|$ that for every tree $T$, vertex $t$ of $T$, and closed walk $w_0, w_1, \dots, w_m$ of $T$ with $w_0 = w_m = t$ that passes each edge of $T$ twice, there exist $h \ge 1$ and disjoint non-empty sets $X(v) \subseteq V(B_h)$ $(v \in V(T))$, each inducing a connected subgraph, such that the root of $B_h$ is in $X(t)$, for each edge of $T$ with ends $u, v$ some edge of $B_h$ joins $X(u)$ and $X(v)$, and there are leaves $l_0, \dots, l_m$ of $B_h$, in order from left to right, with $l_k \in X(w_k)$ for $0 \le k \le m$. If $T$ has no edge, then $m = 0$, and we take $h = 1$ and $X(t) = \{l_0\}$, where $l_0$ is the root. Otherwise let $e$ be the first edge of the walk, joining $t$ and $y = w_1$. Deleting $e$ from $T$ leaves two trees, $T_y$ containing $y$ and $T_t$ containing $t$. Since the walk passes $e$ twice and $e$ is the only edge between $T_y$ and $T_t$, the walk consists of $t$, then a closed walk of $T_y$ from $y$, then a closed walk of $T_t$ from $t$, and these two walks pass each edge of $T_y$ and of $T_t$ twice. By induction they have sets as above. Sets in $B_{h'}$ give sets in $B_{h' + 1}$, by copying them into the subtree at the left child of the root and adding the root to $X(t)$, and so we may assume that both are in $B_{h'}$ for the same $h'$. Let $h = h' + 2$, let $a$ and $b$ be the children of the root of $B_h$, and let $a_1, a_2$ and $b_1, b_2$ be the children of $a$ and of $b$, all from left to right. The subtrees at $a_2$ and $b_1$ are copies of $B_{h'}$. Copy the sets for $T_y$ into the subtree at $a_2$ and the sets for $T_t$ into the subtree at $b_1$, and add to $X(t)$ the root, $a$, $b$, and the path from $a_1$ down to its leftmost leaf, which we take as $l_0$. Then $X(t)$ induces a connected subgraph, since it contains $b_1$, and the edge from $a$ to $a_2 \in X(y)$ joins $X(t)$ and $X(y)$. For $l_1, \dots, l_m$ we take the leaves chosen for the walk of $T_y$ and then those chosen for the walk of $T_t$. They are in order from left to right, after $l_0$, since the subtree at $a_2$ lies to the left of the subtree at $b_1$. This completes the induction.

Now apply this to $F$, $t$ and $W$, and let $\kappa \ge 1$ and $\lambda = h + \kappa$. The subtree of $Y_\lambda$ at a neighbour of its centre is a copy of $B_\lambda$, whose leaves are end-vertices of $Y_\lambda$ in the order of the central circle, and its vertices at distance at most $h - 1$ from its root form a copy of $B_h$. Use the sets $X(v)$ in this copy, and add to $X(w_k)$, for $0 \le k \le m$, the descendants of $l_k$ that are not end-vertices of $Y_\lambda$. This gives a model of $F$ as a minor of $Y_\lambda$ with its end-vertices deleted, in which the $2^\kappa$ end-vertices below $l_k$ form an interval of the central circle, \textit{owned} by $k$, whose vertices are adjacent to vertices of $X(w_k)$. These intervals are disjoint and in the order of $W$ around the central circle, and their length $2^\kappa$ is as large as we wish.

\textit{Step 5: the induction.} We show that $G_j$ has a cyclic model for $j = d, d - 1, \dots, 1$. For $j = 1$, since $G_1 = G$, this proves (2.3).

Let $j = d$, let $p = |V(C_d)|$, and let $g_1, \dots, g_q$ be the edges joining $F$ and $C_d$, listed as in Step 3, where $g_i$ leaves $F$ at the visit $k_i$ and ends at $x_i \in V(C_d)$. (If $F$ is null, then $q = 0$.) If $F$ is a tree, take the model of $F$ of Step 4 with $2^\kappa \ge q(p + 1)$. Since $k_1 \le \dots \le k_q$, we can choose for each $i$ an end-vertex $z_i$ in the interval owned by $k_i$, so that $z_1, \dots, z_q$ are in order around the central circle, with at least $p$ other vertices of this circle between any two of them that are consecutive in the cyclic order. Since $x_1, \dots, x_q$ are in the cyclic order of $C_d$, we can then divide the central circle into intervals $I(x)$ $(x \in V(C_d))$, each with at least one vertex and in the clockwise cyclic order of $C_d$, with $z_i \in I(x_i)$ for each $i$. (The spare vertices between consecutive $z_i$ leave room for the intervals of the vertices of $C_d$ that lie strictly between $x_i$ and $x_{i+1}$ in the cyclic order.) If $q = 0$, we only need $3 \cdot 2^{\lambda - 1} \ge p$. Represent each $x \in V(C_d)$ by $I(x)$, each edge of $C_d$ by the edge of the central circle between the intervals of its ends, which are consecutive since $C_d$ is a circuit, each vertex and edge of $F$ as in Step 4, and each $g_i$ by the edge of $Y_\lambda$ joining $z_i$ to its neighbour, which is in the set of the end $w_{k_i}$ of $g_i$ in $F$. This is a cyclic model of $G_d$ in $N(1, \lambda)$.

Now let $j < d$, let $p = |V(C_j)|$, and let $g_1, \dots, g_q$ be the edges between $C_j$ and $C_{j+1}$, listed as in Step 3, where $g_i$ joins $y_i \in V(C_j)$ and $x_i \in V(C_{j+1})$. Take a cyclic model of $G_{j+1}$ in $N(d - j, \lambda)$ whose intervals each have at least $q(p + 1)$ vertices, which exists by the first fact of Step 4. Let $I(x)$ be its interval for $x \in V(C_{j+1})$. The graph $N(d - j + 1, \lambda)$ consists of a copy of $N(d - j, \lambda)$, a new outer circle, and a radial edge joining each vertex of the new circle to the vertex of the old outer circle on the same radial line. Use the cyclic model of $G_{j+1}$ in the copy. Since $x_1, \dots, x_q$ are in the cyclic order of $C_{j+1}$ and the intervals are long, we can choose radial lines $\ell_1, \dots, \ell_q$, in order around the circles, such that $\ell_i$ meets the old outer circle in $I(x_i)$, with at least $p$ other radial lines between any two of them that are consecutive in the cyclic order. Since $y_1, \dots, y_q$ are in the cyclic order of $C_j$, we can then divide the new outer circle into intervals $I'(y)$ $(y \in V(C_j))$, each with at least one vertex and in the clockwise cyclic order of $C_j$, such that $\ell_i$ meets it in $I'(y_i)$, as in the case $j = d$. Represent each $y \in V(C_j)$ by $I'(y)$, each edge of $C_j$ by the edge of the new circle between the intervals of its ends, and each $g_i$ by the radial edge of $\ell_i$ joining the two outer circles. With the model of $G_{j+1}$, this is a cyclic model of $G_j$ in $N(d - j + 1, \lambda)$, by Step 2. This completes the induction and the proof of (2.3).\fixnote{fix:2.3-proof-added}{The paper omits this part of the proof. The argument above is
the one used in the formalization.}}

\stmt
(2.4) \textit{If $G$ is a planar graph with radius $\le d$ then $G$ is isomorphic to a minor of $N(d, \lambda)$ for some sufficiently large $\lambda$.}

This follows from (2.1), (2.2), and (2.3).

\stmt
(2.5) \textit{For all $d \ge 0$, $\lambda \ge 1$, the graph $N(d, \lambda)$ has tree-width $\le 3d + 1$.}

\textit{Proof.} If $d = 0$ the result is clear, and we assume $d > 0$. We construct a tree-decomposition of $N(d, \lambda)$ as follows. Let $T$ be the tree obtained from $N(d, \lambda)$ by deleting the $3 \cdot 2^{\lambda-1} \times d$ cylinder, and let $C$ be the central circle of the cylinder. For each vertex $v$ of $C$, let $P_v$ be the set of all vertices of the cylinder on the same radial line as $v$. Thus $|P_v| = d$. Let $t$ be a vertex of $T$. Now $t$ is on precisely three regions $R_1$, $R_2$, $R_3$ say (taking a drawing $M$ of $N(d, \lambda)$ as in Fig. 3). For $j = 1, 2, 3$, there are two vertices of $C$ on $R_j$ and they are adjacent. Let them be $u_j$, $v_j$, where $v_j$ follows $u_j$ in the anticlockwise orientation of $C$. Define
\begin{equation*}
X_t = P_{v_1} \cup P_{v_2} \cup P_{v_3} \cup \{t, t'\},
\end{equation*}
where $t' = t$ if $t$ is the ``center'' of $T$, and $t'$ is the neighbour of $t$ closer to the center if $\delm{t'}\fix{t}$ is not the center of $T$.\fixnote{fix:2.5-center}{The vertex $t'$ is being defined here, and the condition is on $t$.} It is easy to verify that $(T, \mathcal{X})$ is a tree-
